\section{Vision and approach}

\subsection{Vision}

This research develops new mathematical frameworks for multi-scale topological analysis of competitive collective systems: bounded collections of agents that coordinate within groups while competing between them. A distinguishing feature is simultaneous operation across multiple spatial scales: individual interactions at the scale of metres, small-group tactical structure at tens of metres, and system-wide organisation beyond that. Standard single-parameter persistent homology applied to such systems conflates topological features from these distinct levels into a single persistence diagram, making scale-specific attribution and tracking of scale dynamics over time impractical. Professional football provides an ideal testbed for developing and validating methods: high-frequency tracking data (10--25~Hz), a well-structured competitive system, a domain-interpretable hierarchy (individual players $\leftrightarrow$ tactical units $\leftrightarrow$ team), and detailed event annotations. This Small Grant develops and validates topological methods on football data at full-season scale; the framework, by design, is transferable to other bounded competitive multi-agent systems with domain-validated interaction scales.

\paragraph{Prior work and the research gap.}
\textit{Collective spatial systems across domains.} The mathematical question of how multi-scale spatial organisation emerges from local interactions recurs across applied mathematics. In physics, active matter self-organises across particle, cluster, and system-wide scales, producing ring-like vortical structures at scale-dependent interaction radii \citep{Marchetti2013, Vicsek2012}. In ecology, collective animal groups, from starling murmurations to fish shoals, coordinate through topological neighbourhood rules rather than metric distance alone \citep{Ballerini2008}. In network science, clique and cavity structure encodes higher-order coordination beyond pairwise edges \citep{Giusti2016}.

\textit{From models of motion to measurement of shape.} These systems are typically studied through agent-based and active-matter models that generate emergent structure, and summarised empirically by geometric and kinematic descriptors: positions, velocities, densities, and dispersion measures, including the team length, width, and centroid measures standard in football analytics \citep{Folgado2014, Gudmundsson2017, FernandezBornn2018}. Such descriptors quantify where agents are and how they move, but not the two-dimensional enclosure (loops and holes) that scale-dependent coordination produces. What is missing is a measurement framework for that structure, not another generative model of it.

\textit{Topology as the measurement tool.} Persistent homology supplies exactly such a framework: it quantifies loops and enclosure rigorously and stably across scale, and is now mature in protein and materials science \citep{Xia2014, Hiraoka2016}, biological aggregation \citep{Topaz2015}, and collective cell motion \citep{Bhaskar2019}.

\textit{The gap, and why football resolves it.} Applied to multi-agent systems, however, persistent homology has so far operated at a single scale, in cooperative or quasi-static settings, or on simulated and small samples \citep{Gu2022, SchindlerBarahona2023}; multiparameter persistence offers simultaneous multi-scale analysis in principle but faces computational barriers at the point-cloud scales relevant here \citep{BotnanLesnick2022}. None combines domain-validated scale decomposition with adaptive Vietoris--Rips filtration on \emph{competitive}, \emph{hierarchical}, \emph{high-frequency} tracking data, validated at population scale. Professional football is the testbed where this becomes achievable: a competitive system, with a domain-interpretable hierarchy, sampled at high frequency and annotated densely enough for population-scale validation.
\textit{What already exists.} The ArXiv preprint \citep{Brown2026} establishes a \emph{constructive} multi-scale pipeline in computational topology: at each time frame and each validated cutoff $\delta$, single-linkage clustering produces a reduced centroid cloud $\tilde P_\delta(t)$; persistent homology is computed on $\tilde P_\delta(t)$ with an adaptive maximum filtration $\varepsilon_{\max}$ that depends on inter-centroid geometry and on $\delta$, so that $H_1$ (loops) is not systematically missed when the cloud is coarse. The output is a \emph{family} of barcode summaries indexed by organisational level, avoiding the scale conflation inherent in a single Vietoris--Rips diagram on all agents. The work supplies explicit geometric cycle representatives for $H_1$ features and statistically links persistence to match events. Ten-match validation shows these regimes are reproducible across independent competitions.
\textit{What this proposal adds.} The preprint establishes \emph{that} scale-specific barcodes exist and are reproducible; this proposal asks \emph{what they tell us at scale}, extending from a 10-match proof of concept to (i)~\emph{distributional} questions---population-level laws of variation for $H_0$/$H_1$ barcodes and for their summaries over hundreds of matches; (ii)~\emph{comparison geometry} in persistence space---distances between diagrams or between persistence landscapes \citep{Bubenik2015, Chazal2014}---to treat tactical systems as distinguishable subsets of barcode or landscape space; (iii)~\emph{functional dynamics}---time-indexed landscapes to characterise stability regimes and transitions within matches. Transfer beyond football follows from the framework's design: any bounded competitive multi-agent system with a domain-interpretable hierarchy of interaction scales admits the same clustering--adaptive filtration--homology construction once cutoffs are validated; football provides the data density and event annotations to calibrate the pipeline. None of (i)--(iii) is demonstrated at championship scale in the preprint.

\paragraph{Mathematical perspective: empirical topology, dynamics, and joint analysis.}
The mathematical contribution goes beyond pipeline generalisation: a full Championship season yields an \emph{empirical distribution} of topological invariants at a scale more typical of physics or materials science than applied topology --- $>500$ matches support estimation of means, variances, and correlations of $H_0$/$H_1$ counts, barcode lengths, and landscape norms at each analysis scale, and hypothesis tests against match-level covariates. \emph{Dynamics} are central: a match is not a single diagram but a path $t \mapsto \mathrm{PD}_\delta(t)$ (or $t \mapsto \lambda_\delta(t)$ in landscape space), so the object of study includes continuity, jumps, and regime shifts in time---not only static fingerprints of ``which formation.'' Coupling those paths to event streams and to established tactical descriptors (possession structure, pressing, phase-of-play labels) \citep{Memmert2017, Gudmundsson2017} tests whether barcode statistics are \emph{orthogonal} to conventional style metrics or recover complementary structure. Tactical fingerprinting is one output of (i)--(ii); the broader mathematical target is a statistically grounded map from competitive dynamics into persistence space.

\paragraph{Excellence and importance.}
Topological Data Analysis (TDA) has proven effective for characterising the shape of complex data, but existing applications examine systems at single scales, in non-competitive scenarios, or with static data. A standard Vietoris--Rips filtration on a multi-agent point cloud conflates topological features from different organisational levels into a single persistence diagram. Our framework addresses this through two methodological contributions: domain-informed hierarchical clustering to decompose point clouds by organisational level, and adaptive filtration ensuring consistent $H_1$ (loop) detection across all scales, a practical necessity since a fixed Vietoris--Rips threshold appropriate at one level produces null results at another. These contributions advance computational topology directly: practical, validated methods for multi-scale persistent homology on hierarchical, competitive, time-evolving point clouds.

Validation across 10 professional football matches (1{,}500 analysis frames, SkillCorner open broadcast tracking throughout, primary demonstrator match 1996435, Sydney~FC vs Adelaide~United) confirms three robust $H_0$ analysis regimes, individual (2.98~m), tactical (12.0~m), and team (30.0~m), with cross-epoch stability scores of 0.96, 0.84, and 1.00 respectively, and two $H_1$ regimes capturing complementary structural information. Individual-scale $H_1$ reveals frequent, transient loops ($97.0\%\pm1.5\%$ presence rate across 10 matches; $95.3\%$ in the primary match, 143/150~frames), whilst tactical-scale $H_1$ reveals rarer but geometrically persistent loops ($19.3\%\pm7.2\%$ presence rate across 10 matches; $12.7\%$ in the primary match, 19/150~frames). Scale non-redundancy is quantified: Spearman $\rho=0.264$ (95\% bootstrap CI $[0.200,\,0.314]$), confirming that neither scale subsumes the other. Real event correlation across 10 matches (104{,}722 event--topology pairs) demonstrates that topological features respond coherently to match dynamics---disruption events are associated with decreased persistence, organisational events with increased persistence ($p<0.001$ for both). Temporal persistence dynamics are match-specific rather than universal, consistent with tactical diversity across matches; the framework provides systematic tools to quantify this variation.

\paragraph{Research objectives.}
Building on this validated foundation, the project pursues three research objectives over 12 months, with a fourth strategic output: (1)~Scale the analysis pipeline to a full Championship season ($\sim$540 matches), stress-testing the framework and establishing population-level topological statistics; (2)~Characterise pattern formation by developing baseline topological fingerprints for different tactical systems, quantifying formation times and distinguishing between tactical configurations in persistence space; (3)~Develop persistence landscape methods \citep{Bubenik2015} for temporal dynamics, classifying stability regimes and characterising transitions within and across matches; (4)~Synthesise full-season results into a Standard Grant application extending the framework to broader competitive systems.

\paragraph{Statistical power and sample-size justification.}
The $\sim$540-match target is chosen so that each objective is adequately powered after accounting for covariate stratification and multiple testing. For Objective~1, the 10-match pilot \citep{Brown2026} gives an across-match standard deviation of $0.072$ for the tactical-scale $H_1$ presence rate (matched by an empirical match-level bootstrap half-width of $0.042$ at $n=10$); a target 95\% confidence-interval half-width of $0.025$, appropriate for subgroup inference by phase-of-play, opponent strength, and venue, requires $n \approx 32$ matches per subgroup, so a full season ($\sim$540 matches) comfortably supports the planned ${\sim}5$--$10$ covariate-stratified subgroups. The binding constraint for temporal contrasts is the same cross-match replication: a pilot linear mixed model of the half-effect on per-window tactical-scale $H_1$ persistence over the 10 matches gives $\hat\beta_1=-0.081$, LMM $p=0.079$, stratified within-match permutation $p=0.051$ \citep{Brown2026}, borderline at $\alpha=0.05$ and consistent with the ``match-specific dynamics'' interpretation preferred here; cross-match replication at full-season scale is what resolves it. For Objective~2, a two-sample comparison in landscape $L^2$-norm between two tactical classes with balanced $n \approx 180$ matches/class detects an effect size of Cohen's $d \ge 0.30$ at $\alpha=0.05$ (power $0.80$) after Benjamini--Hochberg FDR correction \citep{BenjaminiHochberg1995} across the expected $\le 10$ planned contrasts; with three tactical classes this exhausts roughly one-third of the season, leaving capacity for replication folds. The Cohen's $d=0.30$ minimum detectable effect assumes a pooled between-class standard deviation commensurate with the pilot within-frame cross-scale landscape dispersion (Table\,tab:tda\_distances of \citealp{Brown2026}: median $5.67$, IQR $6.95$). This assumption is refined by a pre-registered 20-match labelled pilot at Month~2 before committing to the full Obj~2 run; if the estimated pooled between-class standard deviation moves by more than $25\%$, the matches-per-class target and the FDR contrast budget are updated in the pre-registration amendment (OSF record timestamped) before any confirmatory test is executed, consistent with the standard calibrated pilot practice for pre-registered analyses. For Objective~3, a CUSUM change-point detector on landscape norms with block-bootstrap calibration is powered by within-match frame count ($\approx 4\times 10^{4}$ frames per match at 10\,Hz), so the binding constraint is cross-match replication rather than within-match sample size. A power-analysis vignette is committed to the pre-registered plan (Month~2, \emph{Maximising translation}).

\paragraph{Advancing current understanding.}
The 10-match validation establishes that multi-scale topological structure is a robust feature of competitive spatial systems, not an artefact of any single match. The framework now addresses open questions in applied computational topology: how does one perform meaningful scale-specific persistent homology when the system has known hierarchical organisation? Our answer, clustering-then-filtration with adaptive thresholds, is general, reproducible, and validated on real data. This grant extends this in three directions: (i)~full-season analysis establishes whether topological signatures are stable population features or exhibit meaningful distributional variation, including joint statistics of $H_0$/$H_1$ alongside conventional tactical variables; (ii)~tactical fingerprinting tests whether persistence diagrams can discriminate between tactical systems, connecting TDA to practical classification problems; (iii)~persistence landscape dynamics develops functional representations of temporal evolution---paths in landscape space within and across matches---creating the mathematical apparatus needed for stability regime characterisation and for comparing dynamics, not only static snapshots. Publications target the \emph{Journal of Applied and Computational Topology} and \emph{SIAM Journal on Applied Mathematics}.

\paragraph{Timeliness.}
This research addresses foundational open questions in computational topology for hierarchical, time-evolving systems. Three new theoretical and methodological contributions are developed: (1)~persistence landscape methods for temporal dynamics, characterising stability regimes and regime shifts within and across matches; (2)~comparison geometry in persistence space, enabling topological fingerprinting---discriminating between tactical systems via landscape distances; (3)~empirical distributional theory on topological invariants at championship scale, supporting hypothesis tests against match-level covariates. These contributions advance the mathematical foundations of applied topology directly. The \citet{BotnanLesnick2022} survey signals theoretical maturity in multiparameter persistence; practical validated workflows for domain-informed decomposition remain underdeveloped. A full Championship season now provides, for the first time, a dataset large enough to establish population-level topological statistics. Computational advances in TDA libraries (Ripser, GUDHI) now make full-season multi-scale analysis feasible. Industry partnerships with championship clubs demonstrate real demand for systematic spatial analysis tools. The framework aligns with UKRI priorities for explainable AI, since TDA features are interpretable by construction: each topological feature has a direct geometric realisation.

\paragraph{Impact.}
Direct beneficiaries include the mathematical sciences community (validated methods for multi-scale TDA in competitive systems), sports science researchers (quantitative tools for spatial analysis), and industry partners (championship clubs). Preliminary cross-domain application to armed conflict event data \citep{Brown2026inprep} suggests the three-scale structure transfers across domains with appropriate scale adjustment, defining boundary conditions for when the framework produces meaningful results. The Small Grant provides the full-season evidence base needed to underpin a cross-domain Standard Grant. Indirect beneficiaries include autonomous systems developers (multi-agent coordination), emergency services (crowd management), and conservation organisations (ecosystem monitoring).

\paragraph{Future directions and community use.}
The research space opened by multi-scale persistence on competitive systems is substantially larger than a 12-month Small Grant can cover: inferential statistics on barcode and landscape distributions, systematic comparison to a wider menu of tactical KPIs, landscape-based change detection across leagues, and cross-domain transfer each merit programmes in their own right. This proposal therefore \emph{focuses} on a bounded set of objectives that establish a reproducible full-season baseline, documented methodology, and clearly stated open questions. That combination is designed to lower the barrier for other groups to adopt or extend the pipeline and for the team to pursue a Standard Grant without over-claiming results here. Broader possibilities are summarised in the project outputs (reports, code release, papers) as a roadmap rather than as commitments within this award.

\paragraph{Pathway to larger research programme.}
Building on the validated framework, this project pursues three research objectives over 12 months, with a fourth strategic output. The Standard Grant will extend the framework to broader competitive systems---autonomous coordination, ecological dynamics, crowd monitoring---leveraging the validated methods and full-season evidence base from this Small Grant as foundation. Standard Grant preparation begins in Month~9, incorporating full-season results for submission within 12 months of project completion.

\subsection{Approach}

\paragraph{Project structure and team.}
This 12-month project involves: (i)~the Principal Investigator (PI, 0.2~FTE), responsible for framework development, analysis oversight, and publication; (ii)~two Co-Investigators (combined 0.25~FTE) contributing statistical expertise and domain knowledge; (iii)~a Research Associate (1.0~FTE, 9 months, Months~2--10), responsible for data pipeline implementation, full-season analysis, and landscape methods development; (iv)~partnership with Swansea City~AFC and StatsBomb, providing secured access to Championship tracking data and formation labels. The project is structured around four integrated Objectives (three research objectives plus one strategic output).

\paragraph{Delivering the objectives.}
Three research objectives with clear success criteria, plus a strategic fourth output:

\textit{Objective 1, Full-Season Analysis (Months 1--7).} Scale the validated pipeline to the full Championship season ($\sim$540 matches). PI establishes the concurrent processing pipeline on Supercomputing Wales (Months 1--2) and analyses an initial cohort of 15--20 matches (Months 1--3). RA executes full-season analysis of remaining matches (Months 3--7). Success criteria: ${>}500$ matches processed; population-level $H_0$ and $H_1$ distributions established; cross-match variance characterised; exploratory joint analysis with match-level tactical descriptors (e.g.\ style and phase summaries from partner data) to position barcode statistics relative to established KPIs.

\textit{Objective 2, Tactical Fingerprinting (Months 4--8).} Develop baseline topological fingerprints for different tactical systems. Matches are labelled by \emph{dominant formation}, defined as the modal formation over the first 30 minutes of open play (excluding set pieces and stoppages). Formation labels for the full 2024/25 Championship season are drawn from the commercial StatsBomb feed secured through the Swansea City AFC--StatsBomb partnership, with SkillCorner broadcast annotation as an independent parallel source; inter-rater reliability between the two label streams (Cohen's $\kappa$) is computed on the full overlap and reported before any downstream test, and matches with \(\kappa\)-flagged disagreement are either reconciled by manual review against the broadcast clip or excluded with documented reason (the exclusion list is frozen in the Month-2 pre-registration). If either feed falls below 95\,\% match coverage, the remaining feed is used alone with an added \(\pm 5\,\%\) label-noise sensitivity analysis in the pre-registration update. Persistence diagrams are compared using two complementary metrics: a persistence-image kernel \citep{Adams2017} for vectorised classification, and a landscape $L^2$ distance \citep{Bubenik2015} for continuous-valued contrasts. Success criteria: (i)~topological signatures for $\geq 3$ tactical configurations (e.g.\ 4-4-2, 4-3-3, 3-5-2) with Cohen's $\kappa \geq 0.60$ against StatsBomb labels; (ii)~statistical tests for inter-system differences ($p < 0.05$ with Benjamini--Hochberg FDR correction \citep{BenjaminiHochberg1995}); (iii)~\textbf{non-redundancy}: topological fingerprints must add information beyond standard geometric descriptors (team length, width, surface area; \citealp{Folgado2014}) and pitch-control entropy \citep{FernandezBornn2018}, tested by ANCOVA with partial $\eta^2 \geq 0.05$ for the topological block after adjusting for baseline covariates. The 10-match pilot returns a symmetric proxy (convex-hull area, partial $R^2 = 0.091$ for tactical-scale $H_1$ persistence after controlling for team length, width, and Voronoi dispersion entropy \citep{Brown2026}), supporting feasibility; the ANCOVA power calculation above is independent of this proxy and remains the definitive test at full-season scale.

\textit{Objective 3, Persistence Landscape Dynamics (Months 6--10).} Implement persistence landscape representations \citep{Bubenik2015} for temporal evolution analysis at the individual and tactical scales. \emph{Within-match} paths $t \mapsto \lambda_\delta(t)$ are summarised by functional principal component analysis (FPCA) \citep{RamsaySilverman2005} to extract low-dimensional tactical modes, and regime transitions are detected by a CUSUM change-point procedure \citep{Page1954} on the running $L^2$-norm increment $\| \lambda_\delta(t+\Delta) - \lambda_\delta(t) \|_{L^2}$, with significance assessed against bootstrap confidence bands for persistence landscapes \citep{Chazal2014}. \emph{Between-match} comparison uses Fr\'echet mean landscapes per tactical class and the induced permutation test. Success criteria: functional persistence landscape computation at both analysis scales; reproducible stability-regime segmentation ($\geq 70\%$ frame-level agreement across random sub-samples); change-point alignment with tactical annotation (substitutions, formation changes) at a rate above chance (permutation $p < 0.05$).

\textit{Strategic Output, Standard Grant Preparation (Months 9--12).} Synthesise full-season results, persistence landscape theory, and tactical fingerprinting evidence into a Standard Grant proposal. Identify the key mathematical questions emerging from full-season data that require a larger programme to address.

\paragraph{Feasibility and risk management.}
The technical groundwork is complete. The framework has been validated across 10 matches with computational efficiency (${<}2$\,s per frame on Supercomputing Wales), and the full pipeline, from data loading through persistence computation to statistical analysis, is implemented, tested, and documented in Python. Mature open-source libraries (Ripser, GUDHI, giotto-tda) provide robust infrastructure. Championship broadcast data is secured through the Swansea City AFC--StatsBomb partnership. The primary risk, enhanced GPS data access, is mitigated through multiple independent pathways (Genius Sports UK; Borussia Dortmund, Germany) and the validated broadcast data approach, which demonstrated that framework results are robust without GPS precision. RA recruitment risk is mitigated through early advertising (Month $-2$), competitive salary, and the cross-disciplinary project scope.

\paragraph{Computational budget.}
At ${<}2$\,s per frame (single-threaded, Ripser $H_0{+}H_1$ at the tactical scale), the headline $H_1$ sweep over $\sim$540 matches $\times\,{\sim}45{,}000$ frames requires $\approx 270$ CPU-days; the full programme adds roughly $3\times$ overhead for bootstrap re-sampling (Objective~1), cross-validated landscape distances (Objective~2), and FPCA/CUSUM ablation (Objective~3), giving a conservative estimate of ${\sim}1{,}100$ CPU-days ($\approx 26{,}000$ core-hours). A Supercomputing Wales allocation at this level is well within the PI's current standing access and is embarrassingly parallel by match. Ripser is run in sparse-approximation mode for the tactical-scale filtration when $\varepsilon_{\max}$ induces a dense complex, capping peak memory at ${\approx}4$\,GB/frame.

\paragraph{Risk management (additions).}
Beyond the data-access and recruitment risks above: (i)~\textit{Formation-label noise} --- StatsBomb labels can lag in-match changes. Mitigation: an unsupervised landscape-space clustering run in parallel with the supervised fingerprint analysis; objectives are reported on whichever labelling agrees between routes, with disagreements reported transparently. (ii)~\textit{Scale-drift across leagues} --- the 12.0\,m cutoff was validated on A-League broadcast data; Championship pitches and tracking pipelines differ. Mitigation: a Month-2 gate that re-runs the Section-2.3 validation on the first 20 Championship matches and adjusts $\delta$ before full-season execution if stability scores fall below 0.80. (iii)~\textit{Dual-use} --- individual-player surveillance outputs (re-identification, physical-performance attribution to named players) are explicitly out of scope; ethics protocols restrict released analyses to squad-level aggregates.

\paragraph{Methodology.}
The established multi-scale TDA framework \citep{Brown2026} operates through six stages: (1)~point cloud construction from 22-player positions; (2)~hierarchical clustering at validated cutoff distances (individual 2.98~m, tactical 12.0~m, team 30.0~m); (3)~adaptive filtration using $\varepsilon_{\max} = \max\bigl(P_{75}(d_{ij}), \max(5.0,\, 2\delta)\bigr)$; (4)~persistent homology computation (Ripser) for $H_0$ and $H_1$; (5)~closed cycle identification through geometric realisation of $H_1$ generators; (6)~statistical analysis (Wilcoxon rank-sum, permutation tests, Spearman correlation, FDR correction). Full methodological detail, including sensitivity analysis and ablation studies, is provided in the ArXiv preprint.

New methods developed in this grant: persistence landscape computation from persistence diagrams at each scale; landscape distance metrics for comparing temporal windows; hierarchical clustering in persistence space for tactical system classification; concurrent processing pipeline for full-season execution.

\paragraph{Previous work and building upon it.}
The ArXiv preprint \citep{Brown2026} reports comprehensive 10-match validation: three $H_0$ analysis regimes with cross-epoch stability scores of 0.96, 0.84, and 1.00 (individual, tactical, team); 4{,}515 $H_1$ loops detected across 10 matches at two complementary scales; real event correlation (104{,}722 pairs across 10 matches) demonstrating coherent topological response to match dynamics; sensitivity analysis confirming robust $H_1$ detection across a broad cutoff range (6--14~m); and linkage method comparison (single, complete, Ward) characterising chaining effects. Temporal dynamics are match-specific, consistent with the expectation that tactical context drives persistence evolution. This grant builds on the validated framework: from 10 matches to a full season (scaling), from detection to classification (tactical fingerprinting), and from time-series statistics to functional representations (persistence landscapes).

\paragraph{Maximising translation.}
Paper~1 (methodology, \citealp{Brown2026}) is on ArXiv, with journal submission to JACT during Month~1. Paper~3 (full-season results and persistence landscapes) targets submission in Month~11, with JACT as primary target. A cross-domain application paper \citep{Brown2026inprep} targets submission by Month~4 and provides the Standard Grant's cross-domain evidence strand. Conference presentations at SIAM Dynamical Systems and Applied Topology meetings. Open-source code release (containerised, Apptainer/Docker) at Month~12. Industry impact through quarterly meetings with championship club partners. Standard Grant preparation (Months 9--12) incorporates full-season results.

\paragraph{Open-science commitments.}
To raise the methodological bar and to facilitate community uptake: (i)~a \textbf{preregistered analysis plan} for Objectives 2 and 3 is deposited on the Open Science Framework by Month~2, fixing formation labelling, metric choice, statistical tests, and decision thresholds before the full-season run; (ii)~the pipeline is shipped as a \textbf{containerised release} (Apptainer/Singularity for SCW; Docker for general use) with pinned dependency versions by Month~10; (iii)~a \textbf{synthetic 22-agent swarm exemplar} with known $H_1$ structure is included in the open-source release, enabling the method to be taught and reproduced without access to proprietary tracking data. The release is accepted only when the exemplar reproduces individual-scale $H_1$ presence within $\pm 3$~pp of the pilot primary match ($95.3\%$) and tactical-scale $H_1$ presence within $\pm 5$~pp of $12.7\%$ \citep{Brown2026}, verified by an automated pass/fail test shipped with the container.

\paragraph{Research environment.}
The project benefits from Swansea University's research environment with established strengths in computational mathematics, data science, and sports analytics. Supercomputing Wales (2~Petaflops) enables ${<}2$\,s per-frame processing, supporting concurrent full-season analysis. Secure GDPR-compliant data storage (100~TB) supports commercial partnerships. The Research Office provides partnership negotiation support, fast-tracked ethics approval (4-week turnaround), and dedicated Impact Officer. Championship broadcast data is secured through the Swansea City AFC--StatsBomb agreement, with GPS data pathways via Genius Sports and Borussia Dortmund. The project is delivered by a Research Associate (1.0~FTE, 9 months, Months 2--10) supervised by PI Dr Rowan Brown, with Co-Investigators Professor Liam Kilduff (Sport and Exercise Sciences) and Professor Gibin Powathil (Mathematics) providing domain expertise and mathematical rigour respectively.
